\documentclass[11pt,a4paper]{article}
\usepackage[utf8]{inputenc}
\usepackage[T1]{fontenc}
\usepackage[catalan]{babel}
\usepackage{lmodern}
\usepackage[sfdefault,lining]{FiraSans}
\usepackage[a4paper,top=2.3cm,bottom=2.2cm,left=1.8cm,right=1.8cm,headheight=1.3cm,headsep=0.4cm,footskip=1cm]{geometry}
\usepackage{amsmath,amssymb}
\usepackage{icomma}
\usepackage{graphicx}
\usepackage[table]{xcolor}
\usepackage{tikz}
\usepackage{fontawesome5}
\usepackage{ccicons}
\usepackage{fancyhdr}
\usepackage{lastpage}
\usepackage{tasks}
\usepackage{needspace}
\usepackage[most]{tcolorbox}
\usepackage{array}
\usepackage[hidelinks,pdftitle={Fitxa d'exercicis de la unitat 2: Àlgebra. 1r de Batxillerat CCSS},pdfsubject={Matemàtiques Aplicades a les CCSS I},pdfkeywords={polinomis, Ruffini, teorema del residu, factorització, fraccions algebraiques, equacions de primer i segon grau, biquadrades, radicals, racionals, exponencials, logarítmiques}]{hyperref}

\definecolor{navy}{HTML}{1B2748}
\definecolor{accent}{HTML}{E8603A}
\definecolor{teal}{HTML}{1F9E90}
\definecolor{soft}{HTML}{EEF1F8}
\definecolor{muted}{HTML}{5D6685}

\setlength{\parindent}{0pt}
\setlength{\parskip}{3pt}
\settasks{label=\alph*), label-format=\bfseries\color{navy}, label-width=1.3em, item-indent=1.8em, column-sep=0.6em, after-item-skip=5pt, before-skip=3pt, after-skip=4pt}

% ---------- capçalera i peu ----------
\newcommand{\llicencia}{\href{https://creativecommons.org/licenses/by-nc-sa/4.0/deed.ca}{\ccbyncsa}\ \ Llicència Creative Commons Reconeixement-NoComercial-CompartirIgual 4.0 Internacional (CC BY-NC-SA 4.0)}
\pagestyle{fancy}
\fancyhf{}
\fancyhead[L]{\raisebox{-0.35\height}{\includegraphics[height=0.95cm]{logo-gva.png}}}
\fancyhead[C]{\small\color{muted}Fitxa d'exercicis · Unitat 2: Àlgebra · 1r Batxillerat CCSS}
\fancyhead[R]{\raisebox{-0.35\height}{\includegraphics[height=0.95cm]{logo-vermellar.png}}}
\fancyfoot[L]{\scriptsize\color{muted}\llicencia}
\fancyfoot[R]{\small\color{muted}\thepage}
\renewcommand{\headrulewidth}{0.4pt}
\renewcommand{\headrule}{\color{muted!40}\hrule height \headrulewidth}
\fancypagestyle{primera}{\fancyhf{}\renewcommand{\headrulewidth}{0pt}
  \fancyfoot[L]{\scriptsize\color{muted}\llicencia}
  \fancyfoot[R]{\small\color{muted}\thepage}}

% ---------- nivells i exercicis ----------
\newcommand{\estrelles}[1]{{\color{accent}\ifcase#1\or\faStar\faStar[regular]\faStar[regular]\or\faStar\faStar\faStar[regular]\or\faStar\faStar\faStar\fi}}
\newcounter{exer}
\newcommand{\ex}[1]{\par\needspace{4\baselineskip}\medskip\refstepcounter{exer}\noindent{\color{navy}\bfseries\theexer.}\ \ }

\newcommand{\seccio}[2]{\par\needspace{16\baselineskip}\bigskip
  \begin{tcolorbox}[colback=navy,colframe=navy,coltext=white,arc=2mm,boxrule=0pt,left=3mm,right=3mm,top=1.2mm,bottom=1.2mm]
  {\large\bfseries #1\quad #2}
  \end{tcolorbox}\vspace{-2pt}}

\newcommand{\exf}[1]{\par\medskip\refstepcounter{exer}\noindent{\color{navy}\bfseries\theexer.}\ \ }
\providecommand{\lt}{<}\providecommand{\gt}{>}
\newcolumntype{C}[1]{>{\centering\arraybackslash}p{#1}}
% quadre de recordatori (només a la fitxa d'exercicis)
\newcommand{\recorda}[1]{\par\smallskip\begin{tcolorbox}[colback=accent!7,colframe=accent!55,arc=2mm,boxrule=0.5pt,left=2.5mm,right=2.5mm,top=1mm,bottom=1mm]\small{\color{accent}\faLightbulb[regular]}\ \textbf{Recorda:}\ #1\end{tcolorbox}}
\newcommand{\res}[1]{#1}\newcommand{\rest}[1]{#1}

\def\fmneed{11}% que el títol d'una secció no quede sol al final del full
% =====================================================================
% Capçalera FentMates de les fitxes: un color per a cada curs.
% Els colors són els de la pàgina inicial (LEVELS a assets/js/materials.js);
% el logo ve de _eines/logo/pdf/ (python3 _eines/logo/genera.py).
%
% Ús, al final del preàmbul de la unitat:
%   % =====================================================================
% Capçalera FentMates de les fitxes: un color per a cada curs.
% Els colors són els de la pàgina inicial (LEVELS a assets/js/materials.js);
% el logo ve de _eines/logo/pdf/ (python3 _eines/logo/genera.py).
%
% Ús, al final del preàmbul de la unitat:
%   % =====================================================================
% Capçalera FentMates de les fitxes: un color per a cada curs.
% Els colors són els de la pàgina inicial (LEVELS a assets/js/materials.js);
% el logo ve de _eines/logo/pdf/ (python3 _eines/logo/genera.py).
%
% Ús, al final del preàmbul de la unitat:
%   \input{../comu/fentmates.tex}
%   \fmcurs{1eso}
% i, al començament del document:
%   \capcalera{Matemàtiques}{1r ESO}{Fitxa d'exercicis · Unitat 2}{Divisibilitat}
% Opcional abans de \input: \def\fmneed{14} (espai mínim abans de cada secció).
% =====================================================================
\makeatletter
\def\fm@logodir{../../logo/pdf}
\providecommand{\fmneed}{10}
\def\fm@cursnom{}

% clau del curs, color i color del text damunt del color
\newcommand{\fm@curs}[3]{\@namedef{fm@color@#1}{#2}\@namedef{fm@text@#1}{#3}}
\fm@curs{1eso}{E6B84A}{navy}
\fm@curs{2eso}{3B82F6}{white}
\fm@curs{3eso}{FF7A50}{white}
\fm@curs{4eso-a}{A78BFA}{white}
\fm@curs{4eso-b}{6D28D9}{white}
\fm@curs{1batx-mi}{9B1C1C}{white}
\fm@curs{1batx-ccss}{EF6B6B}{white}
\fm@curs{2batx-mii}{1E6B47}{white}
\fm@curs{2batx-ccss}{2BB3A3}{white}

\newcommand{\fmcurs}[1]{%
  \@ifundefined{fm@color@#1}{\PackageError{fentmates}{Curs desconegut: #1}{}}{}%
  \edef\fm@tmp{\noexpand\definecolor{curs}{HTML}{\@nameuse{fm@color@#1}}%
               \noexpand\colorlet{curstext}{\@nameuse{fm@text@#1}}}\fm@tmp
  \colorlet{cursfosc}{curs!72!black}%
  \def\fm@logo{\fm@logodir/vertical-#1.pdf}%
  \def\fm@marca{\fm@logodir/marca-#1.pdf}}

% ---------- capçalera de la primera pàgina ----------
% #1 matèria (petit), #2 curs, #3 tipus de document i unitat, #4 títol
\newcommand{\capcalera}[4]{%
  \gdef\fm@cursnom{#2}%
  \par\noindent
  \begin{tcolorbox}[colback=curs,colframe=curs,boxrule=0pt,arc=2.2mm,width=3.5cm,height=2.15cm,
      valign=center,left=3mm,right=2mm,top=1mm,bottom=1mm,nobeforeafter,box align=center]
    \color{curstext}{\footnotesize #1}\par\vspace{1pt}{\fontsize{19}{21}\selectfont\bfseries #2\par}%
  \end{tcolorbox}\hspace{5mm}%
  \begin{minipage}[c]{\dimexpr\linewidth-3.5cm-5mm-2.9cm\relax}
    {\color{cursfosc}\small\bfseries\MakeUppercase{#3}}\par\vspace{2pt}
    {\color{navy}\fontsize{20}{23}\selectfont\bfseries #4\par}
  \end{minipage}\hfill
  \begin{minipage}[c]{2.6cm}\raggedleft\includegraphics[width=2.6cm]{\fm@logo}\end{minipage}\par
  \vspace{2pt}}

% ---------- barres de secció del color del curs ----------
\renewcommand{\seccio}[2]{\par\needspace{\fmneed\baselineskip}\bigskip
  \begin{tcolorbox}[colback=curs,colframe=curs,coltext=curstext,arc=2mm,boxrule=0pt,left=3mm,right=3mm,top=1.2mm,bottom=1.2mm]
  {\large\bfseries #1\quad #2}
  \end{tcolorbox}\vspace{-2pt}}

% ---------- estrelles de dificultat del color del curs (a Batxillerat no se'n posen) ----------
\@ifundefined{estrelles}{}{%
  \renewcommand{\estrelles}[1]{{\color{curs}\ifcase#1\or\faStar\faStar[regular]\faStar[regular]\or\faStar\faStar\faStar[regular]\or\faStar\faStar\faStar\fi}}}

% ---------- capçalera i peu de les pàgines ----------
% primera pàgina: logos institucionals al peu
\fancypagestyle{primera}{\fancyhf{}\renewcommand{\headrulewidth}{0pt}
  \fancyfoot[L]{\raisebox{-0.32\height}{\includegraphics[height=0.78cm]{logo-gva.png}}\hspace{5mm}%
    \raisebox{-0.32\height}{\includegraphics[height=0.78cm]{logo-vermellar.png}}}
  \fancyfoot[R]{\scriptsize\color{muted}\href{https://creativecommons.org/licenses/by-nc-sa/4.0/deed.ca}{\ccbyncsa}\ \ CC BY-NC-SA 4.0 · fentmates.org\qquad\small\thepage}}
% la resta: el símbol a l'esquerra i el curs a la dreta, amb el títol al centre com abans
\fancyhead[L]{\raisebox{-0.3\height}{\includegraphics[height=0.62cm]{\fm@marca}}}
\fancyhead[R]{\small\bfseries\color{cursfosc}\fm@cursnom}
\renewcommand{\headrule}{\color{curs!45}\hrule height \headrulewidth}
\makeatother

%   \fmcurs{1eso}
% i, al començament del document:
%   \capcalera{Matemàtiques}{1r ESO}{Fitxa d'exercicis · Unitat 2}{Divisibilitat}
% Opcional abans de \input: \def\fmneed{14} (espai mínim abans de cada secció).
% =====================================================================
\makeatletter
\def\fm@logodir{../../logo/pdf}
\providecommand{\fmneed}{10}
\def\fm@cursnom{}

% clau del curs, color i color del text damunt del color
\newcommand{\fm@curs}[3]{\@namedef{fm@color@#1}{#2}\@namedef{fm@text@#1}{#3}}
\fm@curs{1eso}{E6B84A}{navy}
\fm@curs{2eso}{3B82F6}{white}
\fm@curs{3eso}{FF7A50}{white}
\fm@curs{4eso-a}{A78BFA}{white}
\fm@curs{4eso-b}{6D28D9}{white}
\fm@curs{1batx-mi}{9B1C1C}{white}
\fm@curs{1batx-ccss}{EF6B6B}{white}
\fm@curs{2batx-mii}{1E6B47}{white}
\fm@curs{2batx-ccss}{2BB3A3}{white}

\newcommand{\fmcurs}[1]{%
  \@ifundefined{fm@color@#1}{\PackageError{fentmates}{Curs desconegut: #1}{}}{}%
  \edef\fm@tmp{\noexpand\definecolor{curs}{HTML}{\@nameuse{fm@color@#1}}%
               \noexpand\colorlet{curstext}{\@nameuse{fm@text@#1}}}\fm@tmp
  \colorlet{cursfosc}{curs!72!black}%
  \def\fm@logo{\fm@logodir/vertical-#1.pdf}%
  \def\fm@marca{\fm@logodir/marca-#1.pdf}}

% ---------- capçalera de la primera pàgina ----------
% #1 matèria (petit), #2 curs, #3 tipus de document i unitat, #4 títol
\newcommand{\capcalera}[4]{%
  \gdef\fm@cursnom{#2}%
  \par\noindent
  \begin{tcolorbox}[colback=curs,colframe=curs,boxrule=0pt,arc=2.2mm,width=3.5cm,height=2.15cm,
      valign=center,left=3mm,right=2mm,top=1mm,bottom=1mm,nobeforeafter,box align=center]
    \color{curstext}{\footnotesize #1}\par\vspace{1pt}{\fontsize{19}{21}\selectfont\bfseries #2\par}%
  \end{tcolorbox}\hspace{5mm}%
  \begin{minipage}[c]{\dimexpr\linewidth-3.5cm-5mm-2.9cm\relax}
    {\color{cursfosc}\small\bfseries\MakeUppercase{#3}}\par\vspace{2pt}
    {\color{navy}\fontsize{20}{23}\selectfont\bfseries #4\par}
  \end{minipage}\hfill
  \begin{minipage}[c]{2.6cm}\raggedleft\includegraphics[width=2.6cm]{\fm@logo}\end{minipage}\par
  \vspace{2pt}}

% ---------- barres de secció del color del curs ----------
\renewcommand{\seccio}[2]{\par\needspace{\fmneed\baselineskip}\bigskip
  \begin{tcolorbox}[colback=curs,colframe=curs,coltext=curstext,arc=2mm,boxrule=0pt,left=3mm,right=3mm,top=1.2mm,bottom=1.2mm]
  {\large\bfseries #1\quad #2}
  \end{tcolorbox}\vspace{-2pt}}

% ---------- estrelles de dificultat del color del curs (a Batxillerat no se'n posen) ----------
\@ifundefined{estrelles}{}{%
  \renewcommand{\estrelles}[1]{{\color{curs}\ifcase#1\or\faStar\faStar[regular]\faStar[regular]\or\faStar\faStar\faStar[regular]\or\faStar\faStar\faStar\fi}}}

% ---------- capçalera i peu de les pàgines ----------
% primera pàgina: logos institucionals al peu
\fancypagestyle{primera}{\fancyhf{}\renewcommand{\headrulewidth}{0pt}
  \fancyfoot[L]{\raisebox{-0.32\height}{\includegraphics[height=0.78cm]{logo-gva.png}}\hspace{5mm}%
    \raisebox{-0.32\height}{\includegraphics[height=0.78cm]{logo-vermellar.png}}}
  \fancyfoot[R]{\scriptsize\color{muted}\href{https://creativecommons.org/licenses/by-nc-sa/4.0/deed.ca}{\ccbyncsa}\ \ CC BY-NC-SA 4.0 · fentmates.org\qquad\small\thepage}}
% la resta: el símbol a l'esquerra i el curs a la dreta, amb el títol al centre com abans
\fancyhead[L]{\raisebox{-0.3\height}{\includegraphics[height=0.62cm]{\fm@marca}}}
\fancyhead[R]{\small\bfseries\color{cursfosc}\fm@cursnom}
\renewcommand{\headrule}{\color{curs!45}\hrule height \headrulewidth}
\makeatother

%   \fmcurs{1eso}
% i, al començament del document:
%   \capcalera{Matemàtiques}{1r ESO}{Fitxa d'exercicis · Unitat 2}{Divisibilitat}
% Opcional abans de \input: \def\fmneed{14} (espai mínim abans de cada secció).
% =====================================================================
\makeatletter
\def\fm@logodir{../../logo/pdf}
\providecommand{\fmneed}{10}
\def\fm@cursnom{}

% clau del curs, color i color del text damunt del color
\newcommand{\fm@curs}[3]{\@namedef{fm@color@#1}{#2}\@namedef{fm@text@#1}{#3}}
\fm@curs{1eso}{E6B84A}{navy}
\fm@curs{2eso}{3B82F6}{white}
\fm@curs{3eso}{FF7A50}{white}
\fm@curs{4eso-a}{A78BFA}{white}
\fm@curs{4eso-b}{6D28D9}{white}
\fm@curs{1batx-mi}{9B1C1C}{white}
\fm@curs{1batx-ccss}{EF6B6B}{white}
\fm@curs{2batx-mii}{1E6B47}{white}
\fm@curs{2batx-ccss}{2BB3A3}{white}

\newcommand{\fmcurs}[1]{%
  \@ifundefined{fm@color@#1}{\PackageError{fentmates}{Curs desconegut: #1}{}}{}%
  \edef\fm@tmp{\noexpand\definecolor{curs}{HTML}{\@nameuse{fm@color@#1}}%
               \noexpand\colorlet{curstext}{\@nameuse{fm@text@#1}}}\fm@tmp
  \colorlet{cursfosc}{curs!72!black}%
  \def\fm@logo{\fm@logodir/vertical-#1.pdf}%
  \def\fm@marca{\fm@logodir/marca-#1.pdf}}

% ---------- capçalera de la primera pàgina ----------
% #1 matèria (petit), #2 curs, #3 tipus de document i unitat, #4 títol
\newcommand{\capcalera}[4]{%
  \gdef\fm@cursnom{#2}%
  \par\noindent
  \begin{tcolorbox}[colback=curs,colframe=curs,boxrule=0pt,arc=2.2mm,width=3.5cm,height=2.15cm,
      valign=center,left=3mm,right=2mm,top=1mm,bottom=1mm,nobeforeafter,box align=center]
    \color{curstext}{\footnotesize #1}\par\vspace{1pt}{\fontsize{19}{21}\selectfont\bfseries #2\par}%
  \end{tcolorbox}\hspace{5mm}%
  \begin{minipage}[c]{\dimexpr\linewidth-3.5cm-5mm-2.9cm\relax}
    {\color{cursfosc}\small\bfseries\MakeUppercase{#3}}\par\vspace{2pt}
    {\color{navy}\fontsize{20}{23}\selectfont\bfseries #4\par}
  \end{minipage}\hfill
  \begin{minipage}[c]{2.6cm}\raggedleft\includegraphics[width=2.6cm]{\fm@logo}\end{minipage}\par
  \vspace{2pt}}

% ---------- barres de secció del color del curs ----------
\renewcommand{\seccio}[2]{\par\needspace{\fmneed\baselineskip}\bigskip
  \begin{tcolorbox}[colback=curs,colframe=curs,coltext=curstext,arc=2mm,boxrule=0pt,left=3mm,right=3mm,top=1.2mm,bottom=1.2mm]
  {\large\bfseries #1\quad #2}
  \end{tcolorbox}\vspace{-2pt}}

% ---------- estrelles de dificultat del color del curs (a Batxillerat no se'n posen) ----------
\@ifundefined{estrelles}{}{%
  \renewcommand{\estrelles}[1]{{\color{curs}\ifcase#1\or\faStar\faStar[regular]\faStar[regular]\or\faStar\faStar\faStar[regular]\or\faStar\faStar\faStar\fi}}}

% ---------- capçalera i peu de les pàgines ----------
% primera pàgina: logos institucionals al peu
\fancypagestyle{primera}{\fancyhf{}\renewcommand{\headrulewidth}{0pt}
  \fancyfoot[L]{\raisebox{-0.32\height}{\includegraphics[height=0.78cm]{logo-gva.png}}\hspace{5mm}%
    \raisebox{-0.32\height}{\includegraphics[height=0.78cm]{logo-vermellar.png}}}
  \fancyfoot[R]{\scriptsize\color{muted}\href{https://creativecommons.org/licenses/by-nc-sa/4.0/deed.ca}{\ccbyncsa}\ \ CC BY-NC-SA 4.0 · fentmates.org\qquad\small\thepage}}
% la resta: el símbol a l'esquerra i el curs a la dreta, amb el títol al centre com abans
\fancyhead[L]{\raisebox{-0.3\height}{\includegraphics[height=0.62cm]{\fm@marca}}}
\fancyhead[R]{\small\bfseries\color{cursfosc}\fm@cursnom}
\renewcommand{\headrule}{\color{curs!45}\hrule height \headrulewidth}
\makeatother

\fmcurs{1batx-ccss}

\begin{document}
\thispagestyle{primera}

\capcalera{Matemàtiques CCSS}{1r Batx}{Fitxa d'exercicis · Unitat 2}{\raggedright Àlgebra}
\vspace{10pt}
Nom i cognoms: \rule{8.4cm}{0.4pt}\hfill Grup: \rule{1.6cm}{0.4pt}\hfill Data: \rule{2.2cm}{0.4pt}
\vspace{8pt}

\begin{tcolorbox}[colback=soft,colframe=soft,arc=2mm,left=3mm,right=3mm,top=1.5mm,bottom=1.5mm]
\small Resol els exercicis al quadern mostrant tots els passos. Dona els resultats de manera exacta i, si cal, també aproximats. En les equacions amb radicals, denominadors amb $x$ o logaritmes, comprova sempre les solucions. En els problemes, indica què és la incògnita i respon amb una frase.
\end{tcolorbox}

% =====================================================
\seccio{1}{Polinomis. Operacions}

\exf{1} Donats els polinomis $P(x)=2x^3-x+3$ i $Q(x)=x^2-2x+1$, calcula:\begin{tasks}(4)
  \task $P(x)+Q(x)$
  \task $P(x)-2Q(x)$
  \task $P(x)\cdot Q(x)$
  \task $\big(Q(x)\big)^2$
\end{tasks}

\ex{1} Simplifica les expressions polinòmiques següents:\begin{tasks}(1)
  \task $2(x-2)^2-3(3x+2)-2(3x-2)(3x+2)$
  \task $(3x+2)^2+2(2x-3)^2-(x-5)(2x-5)$
  \task $(2x^2-2x-13)(3x^2-2x)-3x$
  \task $(x^2+4)(x^2-4)(x-2)+2x(x^3-10)x$
  \task $\left(\frac23x-\frac35\right)\left(\frac32x^2+\frac25\right)+\frac{6}{25}$
\end{tasks}\recorda{$(a\pm b)^2=a^2\pm2ab+b^2$, \ $(a+b)(a-b)=a^2-b^2$. Desenvolupa cada producte i redueix els termes semblants.}

\ex{2} Calcula $(4x^2-3x+2)^2$. (Pots usar que $(a+b+c)^2=a^2+b^2+c^2+2ab+2ac+2bc$.)

\ex{2} Calcula els valors de $a$ i $b$ perquè es verifiquen les igualtats següents:\begin{tasks}(1)
  \task $(3x^2-5x+a)(2x^2+bx)=6x^4-19x^3+19x^2-6x$
  \task $2(x+3)^2-3(x+a)=2x^2+bx+3$
\end{tasks}

\ex{1} Efectua les divisions següents i indica'n el quocient i el residu:\begin{tasks}(1)
  \task $(3x^3+2x^2+x-5):(3x^2+2)$
  \task $(3x^4-2x^2-x+4):(x^2+x+2)$
  \task $(x^3-3x^2-x+6):(2x^3-2x+3)$
  \task $(6x^4-x^2+3x):(2x^2+1)$
\end{tasks}\recorda{Comprova-ho amb la prova de la divisió: $D(x)=d(x)\cdot c(x)+r(x)$, amb el grau de $r(x)$ menor que el de $d(x)$.}


% =====================================================
\seccio{2}{Regla de Ruffini. Teorema del residu}

\exf{1} Efectua les divisions següents aplicant la regla de Ruffini:\begin{tasks}(2)
  \task $(3x^3+2x^2+x-5):(x+2)$
  \task $(3x^4-2x^2-x+4):(x+2)$
  \task $(x^3-3x^2-x+6):(2x+3)$
  \task $(4x^4-3x^2+5):(-x+5)$
  \task $(2x^5-12x^3-6x^2-117):(x-3)$
\end{tasks}\recorda{Si el divisor és $ax+b$, fes Ruffini amb $-\frac ba$ i divideix el quocient entre $a$; el residu no canvia.}

\ex{1} Calcula el valor numèric del polinomi $P(x)=x^2-2x-3$ per a $x=-1$ i $x=3$. Podries trobar un altre valor de $x$ per al qual el valor del polinomi siga zero? Raona la resposta.

\ex{1} Comprova si els nombres següents són arrels del polinomi $P(x)=x^4+3x^3-2x^2+6x-8$:\begin{tasks}(5)
  \task $x=1$
  \task $x=2$
  \task $x=-1$
  \task $x=-4$
  \task $x=-3$
\end{tasks}

\ex{1} Comprova si $M(x)=2x^3-5x^2+4x-4$ és divisible per $x-2$ i, si ho és, troba un polinomi $N(x)$ que permeta escriure $M(x)=(x-2)\cdot N(x)$.

\ex{1} Calcula, sense fer la divisió, el residu de les divisions següents:\begin{tasks}(2)
  \task $(x^3-2x+5):(x+2)$
  \task $(x^{100}-2x^{51}+1):(x+1)$
  \task $(2x^4-3x^2+x-1):(x-1)$
  \task $(x^5-32):(x-2)$
\end{tasks}\recorda{Teorema del residu: el residu de dividir $P(x)$ entre $x-a$ és $P(a)$.}

\ex{2} Calcula el valor de $k$ perquè el polinomi:\begin{tasks}(1)
  \task $P(x)=x^3+x^2-2x+k$ siga divisible per $x-2$
  \task $P(x)=x^3-2x^2+kx+4$ siga divisible per $x+2$
  \task $P(x)=(k-3)x^3-15x+k$ siga divisible per $x+4$
  \task $P(x)=x^3-(k-3)x^2+(k-5)x-3$ siga divisible per $x-3$
\end{tasks}

\ex{2} Utilitza Ruffini (o el teorema del residu) per a calcular el valor de $a$ i $b$ que fa que les divisions següents siguen exactes:\begin{tasks}(1)
  \task $(2x^3+5x^2-2x+a):(x+3)$
  \task $(2x^4+3x^3-2x^2-ax+3):(2x+3)$
  \task $(-6x^4+11x^3+ax^2+5x+4):(-3x+4)$
  \task $(4x^4+6x^3-13x^2+x+b):(-x+5)$
\end{tasks}

\ex{2} Calcula el valor de $k$ perquè la divisió de $x^3-(3k+2)x^2-3$ entre $(x+3)$:\begin{tasks}(2)
  \task Siga exacta.
  \task Tinga residu 15.
\end{tasks}

\ex{2} Troba el valor de $m$ perquè el polinomi $2x^4+9x^3+2x^2-6x+3m$ tinga residu 12 en dividir-lo per $x+2$, sense fer la divisió.


% =====================================================
\seccio{3}{Factorització de polinomis}

\exf{1} Determina les arrels dels polinomis següents:\begin{tasks}(2)
  \task $(x-3)(x+5)(x-2)$
  \task $x(x-2)^3(2x+1)$
  \task $(2x-1)(3x+2)(x+1)^3$
  \task $x^3-3x^2-6x+8$
  \task $x^3+8x^2+17x+10$
  \task $2x^5-11x^4+21x^3-16x^2+4x$
\end{tasks}

\ex{2} Determina, si existeixen, les arrels enteres de cadascun dels polinomis següents i factoritza'ls:\begin{tasks}(2)
  \task $P(x)=x^3-2x^2-5x+6$
  \task $P(x)=x^3+x^2-5x+3$
  \task $P(x)=2x^4-x^3-5x^2+x+3$
  \task $P(x)=4x^3-4x^2-11x+6$
  \task $P(x)=x^4-4x^2$
  \task $P(x)=6x^5+11x^4+3x^3-3x^2-x$
  \task $P(x)=x^6-4x^4-x^2+4$
  \task $P(x)=3x^2-12x+12$
  \task $P(x)=\frac{x^4}{4}+\frac34x^2+\frac{9}{16}$ \ (comprova que no té arrels reals)
\end{tasks}\recorda{Primer, trau factor comú i busca identitats notables. Les arrels enteres són divisors del terme independent: prova-les amb Ruffini fins que quede un polinomi de grau 2.}

\ex{2} D'un polinomi de segon grau, $P(x)$, sabem que $P(1)=-6$, $P(0)=-3$ i que una de les seues arrels és 3. Determina'l i factoritza'l.

\ex{1} Troba el valor de $m$ perquè el polinomi $mx^3-6x^2-4x+8$ tinga 2 com a arrel. Factoritza'l després.

\ex{2} Determina $a$ i $b$ de manera que el polinomi $x^3+ax^2+bx-6$ siga divisible per $x-2$ i per $x+3$.

\ex{2} Troba el polinomi de grau 3 que té les arrels $1$, $-2$ i $3$ i que compleix $P(0)=12$.


% =====================================================
\seccio{4}{Fraccions algebraiques}

\exf{2} Simplifica les fraccions algebraiques següents:\begin{tasks}(2)
  \task $\dfrac{2x^2-8}{x^3+x^2-2x}$
  \task $\dfrac{x^3-x}{x^3+3x^2+2x}$
  \task $\dfrac{x^4-2x^3-3x^2}{x^4-9x^2}$
  \task $\dfrac{x^3-3x^2+3x-1}{x^3-2x^2+x}$
  \task $\dfrac{2x^4-5x^3+x^2+5x-3}{2x^3-x^2-7x+6}$
  \task $\dfrac{x^3-27}{x^2-9}$
  \task $\dfrac{x^3-2x^2-9x+18}{x^3-7x^2+16x-12}$
  \task $\dfrac{x^4-5x^3+9x^2-7x+2}{x^4-6x^3+13x^2-12x+4}$
\end{tasks}\recorda{Factoritza el numerador i el denominador i simplifica només els \emph{factors} comuns: $\dfrac{x+3}{x}\neq3$.}

\ex{2} Efectua les sumes i restes següents reduint prèviament a comú denominador:\begin{tasks}(2)
  \task $\dfrac{1}{x+1}+\dfrac{3}{x^2-1}$
  \task $\dfrac{1}{x^2+x}+\dfrac{3}{x+1}$
  \task $\dfrac{1}{2x}-\dfrac{3}{4x^2-4x}$
  \task $\dfrac{1}{x^2+x}+\dfrac{3}{x^2-x}$
  \task $\dfrac{2x}{x-2}+\dfrac{3x-1}{x+2}-\dfrac{1}{x^2-4}$
  \task $\dfrac{1}{(x-1)^2}+\dfrac{2}{x-1}+\dfrac{1}{x^2-1}$
\end{tasks}\recorda{Factoritza els denominadors: el comú denominador és el m.c.m. (factors comuns i no comuns amb el major exponent).}

\ex{2} Efectua els productes i quocients següents i simplifica el resultat:\begin{tasks}(2)
  \task $\dfrac{3x+3}{12x-12}:\dfrac{(x+1)^2}{x^2-1}$
  \task $\dfrac{x^2+2x-3}{(x-2)^3}\cdot\dfrac{(x-2)^2}{x^2-1}$
  \task $\dfrac{6x+4}{x-2}\cdot\dfrac{4x+1}{2x+4}$
  \task $\dfrac{-6x^2}{x^2+2x+1}\cdot\dfrac{x+1}{3x^3}\cdot\dfrac{x+1}{3x}$
  \task $\dfrac{2x^2-7x+6}{x^2+2x+1}\cdot\dfrac{x^2+6x+5}{2x-3}$
  \task $\dfrac{x^2-4}{x^2+x}:\dfrac{x-2}{x+1}$
\end{tasks}\recorda{Factoritza abans de multiplicar: així podràs simplificar sense desenvolupar els productes.}

\ex{3} Calcula i simplifica el resultat de les operacions combinades següents:\begin{tasks}(2)
  \task $\left(\dfrac3x-\dfrac{2x}{x+1}\right)\cdot\dfrac{x^2+x}{x-1}$
  \task $\dfrac{(x-1)^2}{2}\cdot\dfrac{1}{x^2-1}-\dfrac{3x}{(x+1)^2}$
  \task $\left(\dfrac{2x-1}{x+1}-\dfrac{3x}{x-1}\right)\cdot\dfrac{x^3-x}{-x^2-6x+1}$
  \task $\left(\dfrac{x}{x-1}-\dfrac1x\right):\dfrac{x^2-x+1}{x^2-x}$
  \task $\dfrac{1+\frac{2x+1}{x-1}}{\frac{1}{x-1}+\frac{x}{x+1}}$
  \task $\dfrac{6}{2+x}-\dfrac{4}{2-x}+\dfrac{16}{x^2-4}$
  \task $\dfrac{1-\frac1x}{1-\frac1{x^2}}$
  \task $1+\cfrac{1}{1+\cfrac{1}{1+\frac1x}}$
\end{tasks}\recorda{Primer els parèntesis (i els numeradors i denominadors de les fraccions compostes), després productes i quocients i, al final, sumes i restes.}


% =====================================================
\seccio{5}{Equacions polinòmiques}

\exf{1} Resol les equacions de primer grau següents:\begin{tasks}(2)
  \task $3(2x-1)-2(x+4)=5-(x-2)$
  \task $2(x+3)-3(x-1)=4(2-x)+1$
  \task $5x-2(3-x)=3(x+2)-4$
  \task $4(x-1)-(2x+3)=2(x-3)+1$
  \task $3(x+2)-(x-4)=2(x+5)$
  \task $x-\big[2x-(3-x)\big]=5$
\end{tasks}\recorda{Un signe $-$ davant d'un parèntesi canvia el signe de tots els termes de dins: \mbox{$-(x-2)=-x+2$}.}

\ex{2} Resol les equacions de primer grau amb denominadors següents:\begin{tasks}(2)
  \task $\dfrac{2x-3}{4}+\dfrac x2-\dfrac{3x-1}{5}=2x-1$
  \task $\dfrac{x+3}{4}-\dfrac{2x-1}{6}=\dfrac x3-1$
  \task $\dfrac{x-1}{2}-\dfrac{x+2}{3}=1$
  \task $\dfrac x2-\dfrac{x-1}{3}=\dfrac{x+1}{6}+1$
  \task $\dfrac{2(x+1)}{3}-\dfrac{x-3}{4}=1+\dfrac{x}{12}$
  \task $\dfrac{3x-1}{2}-\dfrac{2(x+1)}{5}=x-\dfrac{x-2}{10}$
\end{tasks}\recorda{Multiplica tots els termes pel m.c.m. dels denominadors. Si davant d'una fracció hi ha un $-$, posa el numerador entre parèntesis.}

\ex{1} Resol les equacions de segon grau següents:\begin{tasks}(2)
  \task $x^2-5x+6=0$
  \task $2x^2+3x-2=0$
  \task $3x^2-12=0$
  \task $x^2+4x+5=0$
  \task $4x^2-4x+1=0$
  \task $5x^2-10x=0$
  \task $2(3x-2)+x(x-1)=-4$
  \task $\dfrac{x^2+1}{2}-\dfrac{2x-3}{4}+\dfrac{x^2}{6}=\dfrac{59}{12}$
  \task $(x+1)^3-(x-1)^3=7$
  \task $(2x-1)^2-3(x+1)=x(x-2)$
\end{tasks}\recorda{$ax^2+bx+c=0\Rightarrow x=\dfrac{-b\pm\sqrt{b^2-4ac}}{2a}$. Les incompletes es resolen més ràpid: traent factor comú ($c=0$) o aïllant $x^2$ ($b=0$).}

\ex{2} Resol les equacions biquadrades següents (o que es resolen amb un canvi de variable):\begin{tasks}(2)
  \task $x^4-17x^2+16=0$
  \task $x^6-9x^3+8=0$
  \task $x^4-8x^2-9=0$
  \task $\dfrac{9}{x^2}=10-x^2$
  \task $2(x+1)^4-8x^3-8(x+3)+8=0$
  \task $4x^4-17x^2+4=0$
\end{tasks}

\ex{2} Resol les equacions polinòmiques de grau superior a 2 següents:\begin{tasks}(2)
  \task $6x^4+13x^3-8x^2-17x+6=0$
  \task $2x^4-x^3-3x-18=0$
  \task $6x^4+5x^3=2x^2+x$
  \task $x^3-7x+6=0$
  \task $x^3-3x^2-9x+27=0$
  \task $x^4-x^3-5x^2-x-6=0$
  \task $x^4-1=0$
  \task $x^4+x^3-9x^2-9x=0$
\end{tasks}\recorda{Passa-ho tot a un membre i factoritza (factor comú, Ruffini…): un producte val zero si algun factor val zero.}

\ex{1} Resol les equacions factoritzades següents:\begin{tasks}(2)
  \task $(x^2-3)(x^2-4)=0$
  \task $(3x+2)(x^2-4x)(x^2-9)=0$
  \task $(x^2-4x+4)(x^2+6x+9)=0$
  \task $(x-2)^2(x+1)^3(2x-7)=0$
\end{tasks}


% =====================================================
\seccio{6}{Equacions amb radicals}

\exf{2} Resol les equacions amb radicals següents:\begin{tasks}(2)
  \task $x-\sqrt x=6$
  \task $-\sqrt{2x-3}+1=x$
  \task $x+\sqrt{2x^2+2x-3}=-1$
  \task $\sqrt{\dfrac{7x+1}{4}}=\dfrac{5x-7}{6}$
  \task $\dfrac{\sqrt{2x-1}}{4}=\dfrac{3}{\sqrt{2x-1}}$
  \task $\sqrt{2-5x}+x\sqrt3=0$
  \task $\sqrt{2+\sqrt{x-4}}=\sqrt{12-x}$
  \task $\sqrt x+\sqrt{x-4}=2$
  \task $\sqrt{x+5}+\sqrt{2x+8}=7$
  \task $\sqrt{x+4}+\sqrt{x-1}=5$
  \task $\sqrt{x+3}-\sqrt{x-2}=1$
  \task $\sqrt{1-4x}-2\sqrt{3+x}=\sqrt{3+x}$
\end{tasks}\recorda{Aïlla un radical, eleva al quadrat i repeteix-ho si cal. Comprova \emph{sempre} les solucions en l'equació inicial: en elevar al quadrat poden aparéixer solucions falses.}

\ex{3} Resol les equacions amb radicals no quadrats següents:\begin{tasks}(3)
  \task $\sqrt[3]{x^2-1}+1=x$
  \task $\sqrt[3]{4x-1}=x-4$
  \task $\sqrt{2x+2}-\sqrt[4]{6x+10}=0$
\end{tasks}\recorda{Eleva els dos membres a l'índex de l'arrel (o al m.c.m. dels índexs).}

\ex{2} Calcula un nombre tal que, si se li suma una unitat i després se n'extrau l'arrel quadrada, s'obté el doble que si se li resten 11 unitats i se n'extrau l'arrel quadrada.


% =====================================================
\seccio{7}{Equacions racionals}

\exf{2} Resol les equacions racionals següents:\begin{tasks}(2)
  \task $\dfrac1x+\dfrac1{x^2}+\dfrac1{x^3}=\dfrac78$
  \task $\dfrac{2x}{3}+\dfrac{2x+3}{x-1}=\dfrac{11}{3x-3}$
  \task $\dfrac{2x}{x-2}+\dfrac{3x}{x+2}=\dfrac{6x^2}{x^2-4}$
  \task $\dfrac{3}{x^2-6x+8}+\dfrac{x-5}{x^2-5x+4}=\dfrac{1}{x-4}$
  \task $\dfrac{x}{\sqrt2}+\dfrac{\sqrt2}{x}=\sqrt2\,x$
  \task $\dfrac{x}{x-2}+\dfrac{2x}{2-x}=-x$
  \task $\dfrac{x}{x+1}+\dfrac{x+1}{x}=\dfrac52$
  \task $\dfrac{1+\frac{x+1}{x-1}}{2-\frac{x-1}{x+1}}=2$
  \task $\dfrac{1}{1-\frac{1}{x+1}}=\dfrac{x+1}{x}$
\end{tasks}\recorda{Multiplica pel m.c.m. dels denominadors i, al final, descarta les solucions que anul·len algun denominador.}

\ex{2} Resol $\dfrac{x-2}{x+2}-\dfrac{1}{x-2}=\dfrac{2x-6}{x^2-4}+1$.

\ex{3} Un grup de persones lloga un autobús per 600\,€ a pagar a parts iguals; al final, 10 persones no poden anar al viatge i cadascuna de les que hi van ha de pagar 5\,€ més. Quantes persones hi anaven inicialment?


% =====================================================
\seccio{8}{Equacions exponencials}

\exf{2} Resol les equacions exponencials següents:\begin{tasks}(2)
  \task $4^{x^2+1}=2^{5x+5}$
  \task $4^{(x-2)^2}=262144$
  \task $3^{x+4}=5$
  \task $\dfrac{1}{2^x}=16^{\frac{x(x-1)}{2}}$
  \task $13^{x^2+2x}-\dfrac1{13}=0$
  \task $2^{2x^2-4x-2}=16$
  \task $2\cdot6^{4x-3}=24$
  \task $3\cdot2^x=2\cdot3^x$
  \task $\sqrt[x]{27}=3^{x+2}$
  \task $2^{x+4}-8^x=0$
  \task $7^{x^2-5x+6}=1$
  \task $49^{\frac{x-5}{3}}=\dfrac{1}{\sqrt{7^{x-1}}}$
\end{tasks}\recorda{Si pots, escriu els dos membres amb la mateixa base i iguala els exponents; si no, pren logaritmes.}

\ex{2} Resol les equacions exponencials següents:\begin{tasks}(2)
  \task $2^{x+1}+2^x+2^{x-1}=28$
  \task $5^{x+3}-5^{x-1}-3120=0$
  \task $4^x+4^{x-1}+4^{x-2}=336$
  \task $3^{x+2}+3^{x+1}+3^x+3^{x-1}=120$
  \task $2^{x-1}+2^{x-2}+2^{x-3}+2^{x-4}=960$
  \task $6^x-9\cdot6^{-x}+8=0$
\end{tasks}\recorda{Descompon les potències ($a^{x+1}=a\cdot a^x$, $a^{x-1}=\frac{a^x}{a}$) i trau factor comú $a^x$.}

\ex{3} Resol les equacions exponencials següents fent un canvi de variable:\begin{tasks}(2)
  \task $9^x+5\cdot3^x-24=0$
  \task $2^{2x+1}-5\cdot2^x+2=0$
  \task $3^{2x}-3^{x-1}=3^{x+1}-1$
  \task $2\cdot10^{2x+4}+3\cdot10^{x+2}-5=0$
  \task $3^{x+2}+9^{x-1}=90$
  \task $2^{2x-4}-5\cdot2^{x-3}+1=0$
  \task $4^{x+1}+2^{x+3}=320$
\end{tasks}\recorda{Si hi ha $a^{2x}$ i $a^x$, fes el canvi $t=a^x$; les solucions $t\le0$ no valen.}

\ex{2} Ingressem 5000\,€ en un compte al 4\,\% d'interés compost anual. Quants anys complets han de passar perquè hi haja almenys 8000\,€?

\ex{2} Quant de temps tarda a duplicar-se un capital col·locat al 3\,\% d'interés compost anual? I si l'interés és del 2,5\,\%, quant tarda a augmentar un 20\,\%?


% =====================================================
\seccio{9}{Equacions logarítmiques}

\exf{2} Resol les equacions logarítmiques següents:\begin{tasks}(2)
  \task $\log2+\log(x+3)=\log(x+5)$
  \task $\log(15-2x)=2\log x$
  \task $\log(2x+3)-\log(x-1)=2\log2+2\log3$
  \task $\frac12\log(x^2-3x+2)=\log x$
  \task $\log\left(\dfrac{2x-2}{x}\right)=2\log(x-1)-\log x$
  \task $2\log x=\log\dfrac x2-1$
  \task $\log(4-5x)+\log(2x)=\log(2x-x^2)+\log3$
  \task $\log_2(x+1)+\log_2(x-1)=3$
  \task $\log_2(5x^2+15x+10)-\log_2(x+2)=2$
  \task $\dfrac{\log(4-x)}{\log(x+2)}=2$
\end{tasks}\recorda{Agrupa els logaritmes amb les propietats fins a arribar a $\log A=\log B$ (i aleshores $A=B$). Comprova que tots els arguments són positius.}

\ex{1} Troba el nombre que compleix que, en afegir al seu logaritme decimal el logaritme decimal de 2, el resultat és 1.

\ex{1} Calcula el valor d'un nombre sabent que el doble del seu logaritme decimal és igual a la suma dels logaritmes decimals de 4 i de 9.

\ex{2} Resol $2\log x-\log(x+6)=\log3$.


% =====================================================
\seccio{10}{Problemes}

\exf{1} Una mare té 38 anys i el seu fill, 10. Quants anys han de passar perquè l'edat de la mare siga el doble que la del fill?

\ex{1} La suma de tres nombres naturals consecutius és 93. Quins són?

\ex{1} Un article costa 60,50\,€ amb l'IVA del 21\,\% inclòs. Quin és el preu sense IVA?

\ex{2} Després d'una rebaixa del 20\,\% i d'una altra del 10\,\% sobre el preu ja rebaixat, una jaqueta costa 36\,€. Quant costava abans de les rebaixes? Quin és el percentatge total de rebaixa?

\ex{2} Un comerciant té un tonell de vi de classe A de 5\,€/L. Hi afig 50 litres de vi de classe B de 2\,€/L i 50 litres d'aigua, i la mescla resulta a 4\,€/L. Quants litres de vi de classe A tenia el tonell?

\ex{1} Tenim 30 litres de vi a 5\,€/L i volem afegir-hi vi de 10\,€/L per a obtindre una mescla que isca a 7\,€/L. Quants litres n'hem d'afegir?

\ex{2} En una botiga es ven te blanc a 18\,€/kg i te verd a 14\,€/kg, i també una mescla dels dos a 16,40\,€/kg. Quina és la composició de la mescla?

\ex{1} En un institut hi ha tres cursos. El 60\,\% de l'alumnat estudia primer, el 30\,\% estudia segon i la resta, que són 15 alumnes, estudia tercer. Quants alumnes hi ha a l'institut?

\ex{1} Una fàbrica exporta el 30\,\% de la producció a Àsia, el 40\,\% a Europa, el 15\,\% a Amèrica i la resta, que suposa una facturació de 3 milions d'euros, a Àfrica. Quina és la facturació total?

\ex{2} Un examen té 20 preguntes. Cada encert suma 5 punts i cada errada en resta 2. Una alumna va traure 58 punts. Quantes preguntes va encertar i quantes en va errar, si no en podia deixar cap sense contestar?

\ex{3} Un botiguer inverteix 125\,€ en la compra d'una partida de pomes. En rebutja 20 quilos per defectuoses i ven la resta per 147\,€, i guanya 0,40\,€ per quilo sobre el preu de compra. Quants quilos en va comprar?

\ex{2} Un solar rectangular fa 10 m més de llarg que d'ample i té una superfície de 1200 m$^2$. Quines en són les dimensions?

\ex{2} Calcula les dimensions d'una finca rectangular si el perímetre fa 140 m i la diagonal, 50 m.

\ex{1} Una bassa tarda 5 hores a omplir-se amb la presa d'aigua habitual i 20 hores amb una mànega. Quant de temps cal si fem servir tots dos mètodes alhora?

\ex{2} Dues persones treballant juntes tarden 10 hores a fer un treball, però una treballa el doble de ràpid que l'altra. Quant tardaria cadascuna per separat?

\ex{2} El nombre de malalts de grip en una ciutat al llarg del passat gener ve donat per $y(t)=100+200\cdot e^{0,2t}$, on $t$ són els dies transcorreguts des de l'1 de gener.\begin{tasks}(1)
  \task Quants malalts hi havia l'1 de gener?
  \task Quin dia el nombre de malalts era de 900?
\end{tasks}

\ex{2} Si inicialment hi ha $n$ bacteris en un cultiu i el nombre es duplica cada hora, després de $t$ hores n'hi ha $N=n\cdot2^t$. Si inicialment n'hi ha 4500, quant de temps tardarà el cultiu a arribar a 50\,000 bacteris?

\vfill
\begin{center}\color{muted}\small\faCheckCircle[regular]\ \ Ànim! Comprova els teus resultats amb les eines interactives dels apunts de la unitat.\end{center}

\end{document}
